\documentclass[12pt, oneside]{article}
\usepackage{amsmath, amssymb, amsthm, amscd}
\usepackage{mathrsfs, mathtools, stmaryrd}
\usepackage{geometry}
\geometry{margin=1in}
\usepackage{hyperref}
\usepackage{enumitem}
\usepackage{tikz-cd}
\usepackage{bbm}
\usepackage{bm}

% Theorem environments
\newtheorem{theorem}{Theorem}[section]
\newtheorem{lemma}[theorem]{Lemma}
\newtheorem{proposition}[theorem]{Proposition}
\newtheorem{corollary}[theorem]{Corollary}
\newtheorem{definition}[theorem]{Definition}
\newtheorem{axiom}[theorem]{Axiom}
\newtheorem{remark}[theorem]{Remark}
\newtheorem{example}[theorem]{Example}
\newtheorem{notation}[theorem]{Notation}
\newtheorem{conjecture}[theorem]{Conjecture}

% Special operators
\DeclareMathOperator{\id}{id}
\DeclareMathOperator{\tr}{tr}
\DeclareMathOperator{\Hol}{Hol}
\DeclareMathOperator{\Spec}{Spec}
\DeclareMathOperator{\Ker}{Ker}
\DeclareMathOperator{\Inv}{Inv}
\DeclareMathOperator{\ad}{ad}
\DeclareMathOperator{\Exp}{Exp}
\DeclareMathOperator{\Hom}{Hom}
\DeclareMathOperator{\Obj}{Obj}
\DeclareMathOperator{\Mor}{Mor}
\DeclareMathOperator{\rank}{rank}
\DeclareMathOperator{\img}{img}
\DeclareMathOperator{\coker}{coker}
\DeclareMathOperator{\spn}{span}
\DeclareMathOperator{\supp}{supp}
\DeclareMathOperator{\res}{res}
\DeclareMathOperator{\ev}{ev}
\DeclareMathOperator{\coev}{coev}
\DeclareMathOperator{\ind}{ind}

% Special symbols
\newcommand{\K}{\mathbf{K}}
\newcommand{\N}{\mathcal{N}}
\newcommand{\A}{\mathcal{A}}
\newcommand{\T}{\mathcal{T}}
\newcommand{\Opr}{\mathcal{O}}
\newcommand{\Sg}{\mathbf{\Sigma}}
\newcommand{\Sh}{\bar{\mathcal{T}}}
\newcommand{\Osh}{\bar{\mathcal{O}}}
\newcommand{\Lambdafam}{\mathbf{\Lambda}}
\newcommand{\wtilde}{\widetilde}
\newcommand{\wh}{\widehat}
\newcommand{\xto}{\xrightarrow}
\newcommand{\lra}{\longrightarrow}
\newcommand{\xra}{\xrightarrow}
\newcommand{\mc}[1]{\mathcal{#1}}
\newcommand{\mf}[1]{\mathfrak{#1}}
\newcommand{\ms}[1]{\mathscr{#1}}
\newcommand{\mb}[1]{\mathbb{#1}}
\newcommand{\bA}{\mathbb{A}}
\newcommand{\bB}{\mathbb{B}}
\newcommand{\bC}{\mathbb{C}}
\newcommand{\bN}{\mathbb{N}}
\newcommand{\bQ}{\mathbb{Q}}
\newcommand{\bR}{\mathbb{R}}
\newcommand{\bZ}{\mathbb{Z}}
\newcommand{\cA}{\mathcal{A}}
\newcommand{\cB}{\mathcal{B}}
\newcommand{\cC}{\mathcal{C}}
\newcommand{\cD}{\mathcal{D}}
\newcommand{\cE}{\mathcal{E}}
\newcommand{\cF}{\mathcal{F}}
\newcommand{\cG}{\mathcal{G}}
\newcommand{\cH}{\mathcal{H}}
\newcommand{\cI}{\mathcal{I}}
\newcommand{\cJ}{\mathcal{J}}
\newcommand{\cK}{\mathcal{K}}
\newcommand{\cL}{\mathcal{L}}
\newcommand{\cM}{\mathcal{M}}
\newcommand{\cN}{\mathcal{N}}
\newcommand{\cO}{\mathcal{O}}
\newcommand{\cP}{\mathcal{P}}
\newcommand{\cQ}{\mathcal{Q}}
\newcommand{\cR}{\mathcal{R}}
\newcommand{\cS}{\mathcal{S}}
\newcommand{\cT}{\mathcal{T}}
\newcommand{\cU}{\mathcal{U}}
\newcommand{\cV}{\mathcal{V}}
\newcommand{\cW}{\mathcal{W}}
\newcommand{\cX}{\mathcal{X}}
\newcommand{\cY}{\mathcal{Y}}
\newcommand{\cZ}{\mathcal{Z}}

% Śūnya and Nāgārjuna symbols
\newcommand{\sunya}{\boldsymbol{\varnothing}}
\newcommand{\catkoi}{\mathsf{C}_{\mathsf{I}}}
\newcommand{\catkoii}{\mathsf{C}_{\mathsf{II}}}
\newcommand{\catkoiii}{\mathsf{C}_{\mathsf{III}}}
\newcommand{\catkoiv}{\mathsf{C}_{\mathsf{IV}}}

% Recognition operator
\newcommand{\R}{\mathcal{R}}

\title{Morphic Calculus: \\ A Foundational Framework via Recognition}
\author{Based on the works of \\ [0.5cm] \Large{śūnya -- cut -- morphisum -- operator}}
\date{\today}

\begin{document}

\maketitle

\begin{abstract}
This work presents a complete foundational framework based on the principle that recognition, not computation, is primitive. All mathematical structures are treated as conditional appearances arising from cuts in an undifferentiated kernel. The framework develops: (1) the recognition operator as primitive, (2) morphic objects stable under transformation, (3) coherence as recognition preservation, (4) calculus with respect to arbitrary clocks/measures, (5) spectral invariants as residues of constrained coherence, and (6) complete operator-theoretic foundations. The framework integrates insights from Pāṇinian grammar, Tamil operator algebra, and Hebrew symbol theory as realizations, never as axioms. A four-layer structure (Śūnya, Cut, Morphisum, Operator) provides complete closure.
\end{abstract}

\tableofcontents

\newpage

\part{FOUNDATIONAL PRINCIPLES}

\section{Motivation and Primitive Notion}

\subsection{Motivation}

Most mathematical formalisms begin with construction, computation, or inference. This work begins instead with recognition.

Recognition is not a process applied to objects; it is the prior condition under which objects, relations, and operations become intelligible.

The purpose of morphic calculus is not to compute new quantities, but to formalize how identity persists under transformation.

\subsection{Primitive Notion}

We introduce recognition as a primitive operation, not reducible to:
\begin{itemize}
    \item equality
    \item equivalence
    \item isomorphism
    \item computation
    \item algorithmic decision
\end{itemize}

Recognition is the act by which two expressions are \emph{seen as the same} without mediation. This act cannot be derived; it is assumed.

\subsection{Distinction from Existing Frameworks}

\begin{itemize}
    \item \textbf{Equality} asserts sameness after construction
    \item \textbf{Isomorphism} asserts sameness via mapping
    \item \textbf{Computation} asserts sameness via procedure
\end{itemize}

Recognition asserts sameness prior to all three. Morphic calculus operates before formal structure stabilizes.

\subsection{Morphic Objects}

A morphic object is not defined by internal composition, but by recognition stability. An object is said to be \emph{morphic} if it remains identifiable under admissible transformations without re-computation.

\subsection{Morphic Transformation}

A transformation is morphic if it preserves recognition. Loss of recognition under transformation defines coherence failure.

\subsection{Scope}

This framework will be used to reinterpret:
\begin{itemize}
    \item computation
    \item differential equations
    \item spectral objects
    \item number-theoretic invariants
\end{itemize}
as secondary phenomena arising from partial or constrained recognition.

\section{Recognition Operator}

\subsection{Definition}

\begin{definition}[Recognition]
Recognition is a primitive operation denoted by
\begin{equation}
\R(\mathbf{A},\mathbf{B})
\end{equation}
which asserts that $\mathbf{A}$ and $\mathbf{B}$ are the same without mediation by construction, mapping, or computation.

Recognition is not a function that produces a value. It is a stability assertion.
\end{definition}

\subsection{Properties}

The recognition operator satisfies:

\begin{enumerate}
    \item \textbf{Idempotence}:
    \begin{equation}
    \R(\mathbf{A},\mathbf{A})
    \end{equation}
    
    \item \textbf{Non-constructivity}: Recognition does not imply the existence of a procedure transforming $\mathbf{A}$ into $\mathbf{B}$.
    
    \item \textbf{Pre-computationality}: Recognition may hold even when no finite computation exists between $\mathbf{A}$ and $\mathbf{B}$.
    
    \item \textbf{Stability Criterion}: A transformation $\mathbf{T}$ is admissible if:
    \begin{equation}
    \R(\mathbf{A},\mathbf{B}) \Rightarrow \R(\mathbf{T}(\mathbf{A}),\mathbf{T}(\mathbf{B}))
    \end{equation}
\end{enumerate}

\section{Coherence}

\subsection{Definition}

\begin{definition}[Coherence]
A system is said to be coherent if recognition is preserved under its evolution.

Formally, let $\mathbf{E}_t$ denote the evolution operator of a system. The system is coherent if:
\begin{equation}
\R(\mathbf{A},\mathbf{B}) \Rightarrow \R(\mathbf{E}_t(\mathbf{A}),\mathbf{E}_t(\mathbf{B})) \quad \forall t
\end{equation}
\end{definition}

\subsection{Loss of Coherence}

Loss of coherence occurs when evolution destroys recognition:
\begin{equation}
\R(\mathbf{A},\mathbf{B}) \text{ but } \neg\R(\mathbf{E}_t(\mathbf{A}),\mathbf{E}_t(\mathbf{B}))
\end{equation}
for some $t$. This loss defines instability.

\subsection{Physical Interpretation}

\begin{itemize}
    \item Stable flow = coherent evolution
    \item Turbulence = coherence breakdown
    \item Noise = incoherent spandrels
\end{itemize}

\subsection{Computational Interpretation}

Persistent computation indicates unstable recognition. Systems that require increasing computational effort over time exhibit coherence loss.

\section{Recognition and Computation}

\subsection{General Principle}

Computation arises when recognition cannot be asserted directly.
\begin{itemize}
    \item Where recognition is immediate, computation is unnecessary.
    \item Where recognition fails, computation is introduced as a surrogate.
\end{itemize}
Thus, computation is not fundamental; it is compensatory.

\subsection{Computational Breakdown}

Let $\mathbf{A}$ and $\mathbf{B}$ be two expressions.
\begin{itemize}
    \item If $\R(\mathbf{A},\mathbf{B})$ holds, no computation is required.
    \item If $\R(\mathbf{A},\mathbf{B})$ fails, computation attempts to restore recognition through finite procedures.
\end{itemize}
Computation therefore operates under recognition deficit.

\subsection{Algorithmic Limitation}

An algorithm succeeds if it restores recognition in finite time. An algorithm fails if recognition is not globally preservable under its operations.

This framework predicts the existence of:
\begin{itemize}
    \item non-computable equivalences
    \item undecidable problems
    \item infinite computational cascades
\end{itemize}
without appealing to paradox.

\newpage

\part{COHERENCE AND EVOLUTION}

\section{Coherence and Evolution Equations}

\subsection{PDEs as Coherence Constraints}

Partial differential equations describe the evolution of quantities under local constraints. In this framework, a PDE encodes the requirement that recognition be preserved locally under infinitesimal evolution.

A PDE does not define motion; it defines admissible coherence.

\subsection{Breakdown of PDE Control}

When recognition is preserved globally, PDE evolution remains stable. When recognition cannot be preserved under infinitesimal propagation, the system exhibits:
\begin{itemize}
    \item blow-up
    \item turbulence
    \item singularities
\end{itemize}
These phenomena signal coherence failure, not merely analytical difficulty.

\subsection{Navier--Stokes Reinterpretation}

The Navier--Stokes equations can be interpreted as constraints enforcing coherence in fluid evolution. Turbulence corresponds to regimes where recognition fails to propagate consistently across scales.

This framework reframes the Navier--Stokes problem as a global coherence preservation problem under local evolution rules. (No solution is claimed here.)

\subsection{Universality of Coherence Principle}

The same coherence principle applies to:
\begin{itemize}
    \item reaction--diffusion systems
    \item nonlinear Schrödinger equations
    \item dissipative systems
    \item stochastic PDEs
\end{itemize}
The mathematical difficulty arises from the same source: recognition loss under evolution.

\section{Coherence and Spectral Formation}

\subsection{Spectral Objects as Residues}

Spectral objects arise when coherence is preserved only under restricted transformations. When recognition cannot be maintained globally, it stabilizes along invariant modes.

These invariant modes manifest as:
\begin{itemize}
    \item eigenvalues
    \item spectra
    \item resonances
    \item zeros
\end{itemize}
A spectrum is therefore not primary data; it is a residue of constrained coherence.

\subsection{Spectral Stability}

Let $\mathbf{T}$ be an operator acting on a system. A spectral value corresponds to a mode for which recognition remains invariant under $\mathbf{T}$. Spectral instability indicates partial recognition failure.

\subsection{Interpretation of Zeros and Eigenvalues}

Zeros, poles, and eigenvalues indicate points of maximal tension between:
\begin{itemize}
    \item local coherence
    \item global recognition
\end{itemize}
They are not anomalies; they are structural signatures.

\subsection{Relevance to Number Theory and Geometry}

In number theory and geometry:
\begin{itemize}
    \item spectra encode arithmetic or geometric constraints
    \item zeros mark breakdown of global recognition
\end{itemize}
This framework reinterprets spectral problems as recognition stabilization problems under arithmetic or geometric evolution. (No conjectures are resolved here.)

\subsection{Universality of Spectral Mechanism}

The same mechanism underlies:
\begin{itemize}
    \item spectral geometry
    \item representation theory
    \item quantum spectra
    \item zeta functions
\end{itemize}
All arise from partial preservation of coherence.

\section{Recognition and Representation}

\subsection{Representation Arising from Recognition Deficit}

Representation arises when recognition cannot be maintained directly across transformations. When a system admits multiple admissible transformations that preserve partial recognition, these transformations organize into a group.

Representation theory formalizes how recognition survives under group action.

\subsection{Symmetry as Recognition Invariance}

A symmetry is not a geometric property; it is a recognition invariant. A transformation belongs to a symmetry group if:
\begin{equation}
\R(\mathbf{A},\mathbf{B}) \Rightarrow \R(g\mathbf{A},g\mathbf{B})
\end{equation}
Symmetry therefore encodes what a system cannot forget.

\subsection{Representations as Stabilized Views}

A representation is a stabilized view of recognition under symmetry. Different representations correspond to different modes of partial recognition. Irreducible representations correspond to minimal recognition units.

\subsection{Quantum and Arithmetic Interpretation}

\begin{itemize}
    \item Quantum states arise as representations preserving coherence under evolution.
    \item Arithmetic objects arise as representations preserving recognition under discrete structure.
\end{itemize}
This places quantum mechanics and number theory within the same structural category.

\subsection{Failure Modes}

When no representation preserves recognition globally:
\begin{itemize}
    \item anomalies appear
    \item symmetry breaking occurs
    \item spectra fragment
\end{itemize}
These are not defects; they are structural inevitabilities.

\subsection{Zero as Cut Operator}

In this framework, zero is not treated as a numerical value but as a structural operation that introduces reference by enforcing a distinction. The appearance of zero corresponds to the first admissible comparison, enabling representation at the cost of coherence. All morphic derivatives operate post-cut; irreversibility, curvature, and entropy arise as residues of this operation.

\newpage

\part{CALCULUS WITH RESPECT TO A CHOSEN MEASURE}

\section{Reference--Measure Foundation}

\subsection{The Core Idea}

Classical calculus defines variation with respect to a distinguished coordinate, typically written as $\frac{d}{dx}$. In morphic calculus, no coordinate is privileged a priori. Instead, variation is defined relative to an internal reference structure that measures change in the system under consideration.

The derivative is \textbf{not inherently "d/dx"}. A derivative is always taken \textbf{with respect to something}. What classical calculus hides is this choice.

\subsection{The Fundamental Definition (g-Derivative)}

Let:
\begin{itemize}
    \item $f$ be a quantity of interest
    \item $g$ be any strictly increasing function that we choose as a \textbf{measure / clock / scale}
\end{itemize}

The derivative of $f$ with respect to $g$ is defined as:
\begin{equation}
\frac{df}{dg}(x) := \lim_{h \to 0} \frac{f(x+h) - f(x)}{g(x+h) - g(x)}
\end{equation}
whenever the limit exists.

This is the rate of change of $f$ per unit change of $g$.

\subsection{Reduction to Classical Calculus}

If $g$ is differentiable and $g'(x) \neq 0$, then:
\begin{equation}
\frac{df}{dg}(x) = \frac{f'(x)}{g'(x)}
\end{equation}
Thus classical calculus is the special case $g(x) = x$.

\subsection{Why This Is Not Just a Trick}

Choosing a different $g$ is \textbf{not a change of notation}. It is a \textbf{change of geometry}. When you change the clock:
\begin{itemize}
    \item distances deform
    \item locality can break
    \item memory can appear
    \item symmetries change
    \item conservation laws shift
\end{itemize}

\subsection{Classification by Choice of Clock}

\begin{itemize}
    \item $g(x) = x$ → ordinary calculus
    \item $g = \text{arc length}$ → calculus along curves
    \item $g = \text{potential / action}$ → canonical mechanics
    \item $g = \text{entropy}$ → thermodynamic geometry
    \item $g = \text{scale (RG)}$ → renormalization group flow
    \item $g = \text{probability measure}$ → stochastic calculus
    \item $g = \text{fractal measure}$ → fractional / anomalous diffusion
\end{itemize}

\subsection{Measure-Based Derivative (Radon--Nikodym)}

Instead of choosing a function $g$, we can choose a measure $\mu$. A measure does not have to be smooth; it can be fractal, probabilistic, singular, or history-dependent.

The derivative of $f$ with respect to $\mu$ is defined as the Radon--Nikodym derivative:
\begin{equation}
\frac{df}{d\mu} \quad \text{whenever } f \ll \mu
\end{equation}

This is where:
\begin{itemize}
    \item memory appears
    \item locality disappears
    \item nonlocal dynamics become natural
\end{itemize}

\subsection{The Śūnya Connection}

When the chosen clock becomes flat, collapses distinctions, or loses resolution, the derivative \textbf{vanishes}—not because nothing happens, but because change is no longer distinguishable.

This is the \textbf{return to Śūnya}:
\begin{itemize}
    \item not zero as a number
    \item but zero as \emph{loss of distinguishability}
\end{itemize}
Limits no longer go to points; they go to kernels.

\newpage

\part{MORPHIC CALCULUS: FORMAL DEVELOPMENT}

\section{The Four-Layer Structure}

The framework is organized into four locked layers:

\begin{enumerate}
    \item \textbf{Śūnya layer}: pre-representation (kernel/ground)
    \item \textbf{Cut layer}: distinction operator that creates structure
    \item \textbf{Morphisum layer}: calculus on transformations (morphisms)
    \item \textbf{Operator layer}: transforms + spectra + invariants across representations
\end{enumerate}

\section{Minimal Formal Core}

\subsection{Three Objects}

\begin{enumerate}
    \item \textbf{A state domain} $\mathcal{X}$ — Not necessarily a vector space. Could be a set, manifold, category, grammar-state space, etc.
    
    \item \textbf{A representation functor} $\mathbf{R}$ — A map from "raw state" into a representable space:
    \begin{equation}
    R: \mathcal{X} \rightarrow \mathcal{H}
    \end{equation}
    where $\mathcal{H}$ is typically a Hilbert space, Banach space, function space, or algebra.
    
    \item \textbf{The Śūnya object} $\varnothing_*$ — Not "0". It is the \textbf{kernel of representation}:
    \begin{equation}
    \varnothing_* := \ker(R) \quad \text{(or the maximal indistinguishability class under }R\text{)}
    \end{equation}
    Meaning: everything that cannot be distinguished in the chosen representation.
\end{enumerate}

\subsection{Three Operators}

\begin{enumerate}
    \item \textbf{Cut operator} $\mathcal{C}$ — A cut is an operation that creates distinguishability:
    \begin{equation}
    \mathcal{C}: \varnothing_* \rightsquigarrow \mathcal{X}
    \end{equation}
    Interpretation: emergence of structure from the undifferentiated kernel. Scientific role: phase transition / symmetry breaking / measurement / model selection.
    
    \item \textbf{Morphisum flow} $\Phi_t$ — A dynamics on structure:
    \begin{equation}
    \Phi_t: \mathcal{X} \rightarrow \mathcal{X}, \quad \Phi_{t+s} = \Phi_t \circ \Phi_s
    \end{equation}
    This is the "morphic motion" of forms.
    
    \item \textbf{Return (collapse) operator} $\mathcal{P}$ — A projection back toward the kernel:
    \begin{equation}
    \mathcal{P}: \mathcal{X} \rightarrow \varnothing_*
    \end{equation}
    Scientific role: coarse-graining, thermalization, forgetting, compression, decoherence, convergence. This is the operator that makes "limit" meaningful.
\end{enumerate}

\section{The Calculus Itself}

\subsection{Morphisum Derivative}

Given an observable $f: \mathcal{X} \rightarrow \mathbb{R}$ (or $f: \mathcal{X} \rightarrow \mathcal{H}$), define the derivative along a morphic flow:
\begin{equation}
D_{\Phi}f(x) := \lim_{h \to 0} \frac{f(\Phi_h(x)) - f(x)}{h}
\end{equation}

This is standard-looking, but the power is: $x$ \textbf{can live in a non-coordinate space} (grammar state, symbolic state, categorical object, etc.)—that's the real "new calculus" move.

\subsection{Morphisum Integral}

Accumulation along morphic evolution:
\begin{equation}
\int_0^T f(\Phi_t(x)) \, dt
\end{equation}
Again standard-looking, but now it integrates \textbf{structure change}, not just coordinates.

\subsection{The Śūnya Limit}

Classical limits go to a number. Here we allow "limits that return to kernel":
\begin{equation}
\overset{\varnothing}{\lim_{t \to \infty}} x_t := \mathcal{P}(x_t)
\end{equation}
So the limit is not "a point"; it can be \textbf{collapse into indistinguishability}.

This is exactly what science needs for:
\begin{itemize}
    \item regime change
    \item compression
    \item collapse
    \item renormalization
    \item phase transitions
\end{itemize}

\section{Generalized Derivative via Choice of Measure/Clock}

\subsection{Primitive Objects for Generalized Calculus}

\begin{itemize}
    \item Let $X$ be a state space (not assumed Euclidean).
    \item Let $f: X \rightarrow \mathbb{R}$ (or into a vector space) be an observable.
    \item Let $g: X \rightarrow \mathbb{R}$ be a \textbf{monotone distinguishability function} (called a clock, scale, or measure-generator).
\end{itemize}
No smoothness is assumed unless stated.

\subsection{g-Derivative (Definition)}

The derivative of $f$ with respect to $g$ at $x$ exists if the limit:
\begin{equation}
D_g f(x) = \lim_{\epsilon \to 0} \frac{f(x_\epsilon) - f(x)}{g(x_\epsilon) - g(x)}
\end{equation}
exists for all admissible paths $x_\epsilon \to x$ consistent with the structure on $X$.

This defines \textbf{rate of change of $f$ per unit distinguishability measured by $g$}.

\subsection{Reduction Rule (Consistency)}

If $X = \mathbb{R}$, $g$ is differentiable and $g'(x) \neq 0$, then:
\begin{equation}
D_g f(x) = \frac{f'(x)}{g'(x)}
\end{equation}
Thus classical calculus is the special case $g(x) = x$.

\subsection{Measure Generalization}

Let $\mu$ be a $\sigma$-finite measure on $X$. If $f$ is absolutely continuous with respect to $\mu$, define:
\begin{equation}
D_\mu f := \frac{df}{d\mu}
\end{equation}
(the Radon--Nikodym derivative). This replaces pointwise locality with \textbf{measure locality}.

\subsection{Morphisum Coupling}

Let $\Phi_t: X \rightarrow X$ be a morphic flow. Define the \textbf{morphic g-derivative}:
\begin{equation}
D_{(\Phi,g)} f(x) = \lim_{h \to 0} \frac{f(\Phi_h(x)) - f(x)}{g(\Phi_h(x)) - g(x)}
\end{equation}

This is the fundamental derivative of the new calculus.
\begin{itemize}
    \item Flow gives \textbf{direction of change}
    \item $g$ gives \textbf{measurement of change}
\end{itemize}
Both are independent choices.

\subsection{Collapse / Śūnya Condition}

If for a region $U \subset X$:
\begin{equation}
g(x_1) = g(x_2) \quad \forall x_1, x_2 \in U
\end{equation}
then all derivatives $D_g f$ vanish on $U$, regardless of internal variation of $f$.

This defines \textbf{return to Śūnya}:
\begin{itemize}
    \item not zero value
    \item but zero distinguishability under the chosen measure
\end{itemize}
Limits are kernel-valued, not point-valued.

\subsection{Consequences (Formal)}

When $g$ or $\mu$ is non-Euclidean:
\begin{itemize}
    \item linearity of derivative may fail
    \item locality may fail
    \item memory terms appear naturally
    \item conservation laws depend on invariance of $g$
    \item reciprocity symmetry can break
\end{itemize}
These are \textbf{structural consequences}, not anomalies.

\section{Chain Rule}

\subsection{Setup}

Let:
\begin{itemize}
    \item State space $X$
    \item Observable $f: X \rightarrow Y$
    \item Function $h: Y \rightarrow Z$
    \item Composition observable: $H(x) = (h \circ f)(x)$
    \item Morphic flow $\Phi_t: X \rightarrow X$
    \item Clock $g: X \rightarrow \mathbb{R}$
\end{itemize}

\subsection{Chain Rule (Basic Form)}

If $D_{(\Phi,g)} f(x)$ exists and ordinary derivative $h'(y)$ exists at $y = f(x)$, then:
\begin{equation}
D_{(\Phi,g)}(h \circ f)(x) = h'(f(x)) \cdot D_{(\Phi,g)} f(x)
\end{equation}

\textbf{This is the fundamental chain rule.} No coordinates assumed. No Euclidean metric assumed.

\subsection{Derivation}

Starting from definition:
\begin{equation}
\frac{h(f(\Phi_h(x))) - h(f(x))}{g(\Phi_h(x)) - g(x)} = \frac{h(f(\Phi_h(x))) - h(f(x))}{f(\Phi_h(x)) - f(x)} \cdot \frac{f(\Phi_h(x)) - f(x)}{g(\Phi_h(x)) - g(x)}
\end{equation}
Taking limits yields the result.

\subsection{Chain Rule with Change of Clock}

Let $g_1$ and $g_2$ be two clocks, both monotone along the same flow. Then:
\begin{equation}
D_{(\Phi,g_1)} f = D_{(\Phi,g_2)} f \cdot D_{(\Phi,g_1)} g_2
\end{equation}
This is the \textbf{clock-chain rule}.

\subsection{Measure-Theoretic Chain Rule}

Let $\mu$, $\nu$ be measures on $X$ with $\nu \ll \mu$. Then:
\begin{equation}
\frac{df}{d\nu} = \frac{df}{d\mu} \cdot \frac{d\mu}{d\nu}
\end{equation}
This is the Radon--Nikodym chain rule, now treated as a \textbf{calculus identity}.

\section{Product Rule}

\subsection{Scalar Product Rule}

Assume $f$ and $q$ are $g$-differentiable along $\Phi_t$ at $x$. Then:
\begin{equation}
D_{(\Phi,g)}(fq)(x) = (D_{(\Phi,g)} f)(x)\, q(x) + f(x)\, (D_{(\Phi,g)} q)(x)
\end{equation}

\subsection{Derivation}

\begin{align}
\frac{(fq)(\Phi_h(x)) - (fq)(x)}{g(\Phi_h(x)) - g(x)} &= \frac{[f(\Phi_h(x)) - f(x)]q(x)}{g(\Phi_h(x)) - g(x)} \\
&\quad + \frac{f(\Phi_h(x))[q(\Phi_h(x)) - q(x)]}{g(\Phi_h(x)) - g(x)}
\end{align}
Taking limits and using continuity $f(\Phi_h(x)) \to f(x)$ yields the result.

\subsection{Bilinear Product Rule}

Let $f: X \rightarrow Y$, $q: X \rightarrow W$, and let $\beta: Y \times W \rightarrow Z$ be a continuous bilinear map (e.g., dot product, pairing, tensor contraction). Define $H(x) = \beta(f(x), q(x))$.

If $f, q$ are $g$-differentiable along $\Phi_t$, then:
\begin{equation}
D_{(\Phi,g)} H(x) = \beta(D_{(\Phi,g)} f(x), q(x)) + \beta(f(x), D_{(\Phi,g)} q(x))
\end{equation}

\subsection{Non-commutative Product Rule}

Let $A, B: X \rightarrow \mathcal{A}$ where $\mathcal{A}$ is an associative algebra (operators/matrices). Define pointwise product $(AB)(x) = A(x)B(x)$. If $A, B$ are $g$-differentiable:
\begin{equation}
D_{(\Phi,g)}(AB)(x) = (D_{(\Phi,g)} A)(x) B(x) + A(x) (D_{(\Phi,g)} B)(x)
\end{equation}
No commutativity is assumed.

\subsection{Commutator Derivative}

Define the commutator observable $[A, B](x) = A(x)B(x) - B(x)A(x)$. Then:
\begin{equation}
D_{(\Phi,g)} [A, B](x) = [D_{(\Phi,g)} A, B](x) + [A, D_{(\Phi,g)} B](x)
\end{equation}

\section{Quotient and Inverse Rules}

\subsection{Inverse Rule (Scalar)}

If $f$ is $g$-differentiable at $x$ and $f(x) \neq 0$, then:
\begin{equation}
D_{(\Phi,g)} (f^{-1})(x) = -\frac{D_{(\Phi,g)} f(x)}{f(x)^2}
\end{equation}

\textit{Derivation:} From $f \cdot f^{-1} \equiv 1$, apply product rule:
\begin{equation}
0 = D_{(\Phi,g)}(1) = (D_{(\Phi,g)} f) f^{-1} + f D_{(\Phi,g)}(f^{-1})
\end{equation}
and solve.

\subsection{Quotient Rule (Scalar)}

For $q(x) \neq 0$:
\begin{equation}
D_{(\Phi,g)} \left( \frac{f}{q} \right)(x) = \frac{(D_{(\Phi,g)} f)(x) q(x) - f(x) (D_{(\Phi,g)} q)(x)}{q(x)^2}
\end{equation}

\subsection{Operator Inverse Rule}

Let $A: X \rightarrow \mathcal{A}$ where $\mathcal{A}$ is an associative algebra. Assume $A(x)$ is invertible and $A^{-1}$ is $g$-differentiable. Then:
\begin{equation}
D_{(\Phi,g)} (A^{-1})(x) = -A(x)^{-1} (D_{(\Phi,g)} A)(x) A(x)^{-1}
\end{equation}

\subsection{Operator Quotient Rule}

For operator-valued $A, B$ with $B$ invertible, define $Q(x) = A(x) B(x)^{-1}$. Then:
\begin{equation}
D_{(\Phi,g)} Q = (D_{(\Phi,g)} A) B^{-1} - A B^{-1} (D_{(\Phi,g)} B) B^{-1}
\end{equation}

\section{Integration and Fundamental Theorem}

\subsection{g-Integral (Definition)}

Let $\gamma: [a,b] \rightarrow X$ be a flow-aligned path: $\gamma(t) = \Phi_t(x_0)$. Define the integral of $f$ with respect to $g$ along $\gamma$ as:
\begin{equation}
\int_\gamma f \, dg = \lim_{\max |\Delta g_i| \to 0} \sum_i f(\gamma(t_i)) [g(\gamma(t_{i+1})) - g(\gamma(t_i))]
\end{equation}
This is a Stieltjes-type integral, intrinsic to the calculus.

\subsection{Measure Formulation}

If $\mu$ is a measure on $X$, define:
\begin{equation}
\int f \, d\mu
\end{equation}
as the accumulation of $f$ over distinguishability measured by $\mu$. Clock-integrals are special cases where $d\mu = dg$.

\subsection{Fundamental Theorem (Local Form)}

Let $F: X \rightarrow \mathbb{R}$ satisfy $D_{(\Phi,g)} F(x) = f(x)$ on a flow-connected region where $g$ is strictly monotone. Then for any $x_1, x_2$ on the same flow line:
\begin{equation}
\int_{x_1}^{x_2} f \, dg = F(x_2) - F(x_1)
\end{equation}
This establishes \textbf{integration as the inverse of differentiation} under the chosen clock.

\subsection{Fundamental Theorem (Operator Form)}

Let $D_{(\Phi,g)}$ act as a derivation on observables. Define the integral operator $I_g$ by:
\begin{equation}
(I_g f)(x) := \int_{\gamma_{x_0 \to x}} f \, dg
\end{equation}
Then:
\begin{equation}
D_{(\Phi,g)} \circ I_g = \text{Id}
\end{equation}
on the space of integrable observables modulo constants on flow components.

\subsection{Constants and Śūnya Kernel}

If $f(x) = 0$ almost everywhere with respect to $g$, then $I_g f \in \ker(D_{(\Phi,g)})$. This kernel is the \textbf{Śūnya space of integration constants}:
\begin{itemize}
    \item not numerical constants
    \item equivalence classes of indistinguishable observables under $g$
\end{itemize}

\subsection{Failure Modes}

The fundamental theorem fails precisely when:
\begin{enumerate}
    \item $g$ is not monotone along the path
    \item $g$ collapses distinctions (flat clock)
    \item flow leaves the domain of definition
    \item measure is singular on the path
\end{enumerate}
These correspond to \textbf{regime changes}, not inconsistencies.

\section{Higher-Order Derivatives and Commutators}

\subsection{Second Derivative Along a Single Clocked Flow}

Given flow $\Phi_t$, clock $g$, and first derivative operator $D_{(\Phi,g)}$, define the second derivative by iteration:
\begin{equation}
D_{(\Phi,g)}^2 f := D_{(\Phi,g)} (D_{(\Phi,g)} f)
\end{equation}
Existence requires $D_{(\Phi,g)} f$ exists in a neighborhood along flow and $D_{(\Phi,g)} (D_{(\Phi,g)} f)$ exists.

\subsection{Generator Form (Clock-Normalized)}

Assume a flow parameter $t$ exists and define:
\begin{itemize}
    \item Classical generator along $t$:
    \begin{equation}
    \mathcal{L}_\Phi f(x) := \lim_{h \to 0} \frac{f(\Phi_h(x)) - f(x)}{h}
    \end{equation}
    \item Clock speed (when it exists, nonzero):
    \begin{equation}
    \dot{g}_\Phi(x) := \lim_{h \to 0} \frac{g(\Phi_h(x)) - g(x)}{h}
    \end{equation}
\end{itemize}
Then:
\begin{equation}
D_{(\Phi,g)} f(x) = \frac{1}{\dot{g}_\Phi(x)} \mathcal{L}_\Phi f(x)
\end{equation}
So second derivative expands to:
\begin{equation}
D_{(\Phi,g)}^2 f = D_{(\Phi,g)} \left( \frac{1}{\dot{g}_\Phi} \mathcal{L}_\Phi f \right)
\end{equation}
This makes explicit: \textbf{second order = flow curvature + clock curvature}.

\subsection{Two-Flow Calculus: Mixed Derivatives}

Let there be two flows $\Phi_t$, $\Psi_s$ with clocks $g_\Phi, g_\Psi$ (or one shared clock $g$). Define two derivatives:
\begin{itemize}
    \item $D_\Phi := D_{(\Phi, g_\Phi)}$
    \item $D_\Psi := D_{(\Psi, g_\Psi)}$
\end{itemize}
Mixed derivatives: $D_\Phi D_\Psi f$, $D_\Psi D_\Phi f$.

\subsection{Commutator (Curvature Operator)}

Define the commutator:
\begin{equation}
[D_\Phi, D_\Psi] f := D_\Phi (D_\Psi f) - D_\Psi (D_\Phi f)
\end{equation}
This object measures \textbf{non-integrability / curvature} of the combined evolution.

\subsection{Flatness / Integrability Condition}

The system is (locally) integrable when $[D_\Phi, D_\Psi] = 0$ on the relevant observable algebra. This is the calculus-level definition of "flows commute" under chosen clocks.

\subsection{Lax / Isospectral Evolution}

Operator flow evolution of the form:
\begin{equation}
\frac{d}{dt} L(t) = [A(t), L(t)]
\end{equation}
is an isospectral mechanism. Then:
\begin{itemize}
    \item the spectrum of $L(t)$ is invariant in $t$
    \item spectral data becomes a \textbf{morphic invariant}
\end{itemize}

\newpage

\part{EXPONENTIAL CHANGE OPERATOR FRAMEWORK}

\section{Generators and Exponential Evolution}

\subsection{Generators as Derivations}

Let $\mathcal{A}$ be an observable algebra. A \textbf{change generator} is a derivation $D: \mathcal{A} \rightarrow \mathcal{A}$ satisfying the Leibniz rule:
\begin{equation}
D(fq) = D(f)q + fD(q)
\end{equation}

\subsection{Exponential Evolution Operator}

Define the finite-time change operator by exponential of the generator:
\begin{equation}
E_t := \exp(tD)
\end{equation}
acting on observables. This is the formal internal definition of "finite-time change as exponential of a generator."

\subsection{Semigroup Law}

For all $s, t$ in the admissible parameter set:
\begin{equation}
E_{t+s} = E_t \circ E_s, \quad E_0 = I
\end{equation}
This is the closure condition that turns "change" into a semigroup.

\subsection{Finite-Time Leibniz Law (Multiplicativity)}

For any observables $f, q \in \mathcal{A}$:
\begin{equation}
E_t(fq) = (E_t f)(E_t q)
\end{equation}
This is the "global Leibniz law / finite-time product rule" for pullback-by-flow behavior.

\subsection{Differential Law of the Exponential Flow}

Define $E_t$ to be the unique family satisfying:
\begin{equation}
\frac{d}{dt} E_t f = E_t (Df), \quad E_0 f = f
\end{equation}
This is the internal operational rule: exponential evolution is the flow integrating the generator.

\section{Discrete Change Operator}

\subsection{Exponential Change Operator as Discrete Truncation}

Given a step size $\Delta t$, define the \textbf{discrete change-step operator}:
\begin{equation}
C_{\Delta t} := I + \Delta t \, D
\end{equation}
Then the $n$-step discrete evolution (total time $t = n\Delta t$) is:
\begin{equation}
C_{\Delta t}^n
\end{equation}
This is the canonical "change operator" as first-order (Euler) truncation of the exponential flow.

\subsection{Consistency Condition}

For any observable $f$ in the domain of $D$:
\begin{equation}
\lim_{\Delta t \to 0} \frac{C_{\Delta t} f - f}{\Delta t} = Df
\end{equation}
So $C_{\Delta t}$ is the correct infinitesimal generator in the limit.

\subsection{Convergence Schema (Euler → Exponential)}

Define $t = n\Delta t$, $\Delta t = t/n$. On any subspace where the semigroup is well-defined:
\begin{equation}
\lim_{n \to \infty} \left( I + \frac{t}{n} D \right)^n = \exp(tD) = E_t
\end{equation}
This is the structural convergence statement that restores semigroup closure by exponentiation.

\section{Spectral Modes and Decay Laws}

\subsection{Mode Law (Eigenmode Definition)}

An observable $u \in \mathcal{A}$ is a \textbf{spectral mode} of $D$ with rate $\lambda$ if:
\begin{equation}
Du = \lambda u
\end{equation}
Then:
\begin{itemize}
    \item \textbf{Continuous-time evolution of a mode}:
    \begin{equation}
    E_t u = e^{\lambda t} u
    \end{equation}
    \item \textbf{Discrete-time evolution of a mode under $C_{\Delta t}$}:
    \begin{equation}
    C_{\Delta t} u = (1 + \lambda \Delta t) u
    \end{equation}
    and after $n$ steps:
    \begin{equation}
    C_{\Delta t}^n u = (1 + \lambda \Delta t)^n u
    \end{equation}
\end{itemize}
Mode-by-mode convergence: $(1 + \lambda t/n)^n \to e^{\lambda t}$.

\subsection{Invariants and Kernel}

If $Du = 0$, then $E_t u = u$ for all $t$. So $\ker(D)$ is the invariant subalgebra (zero-modes).

\subsection{Broken-Symmetry Axiom}

If two observables $F, G$ satisfy $C(F) = C(G)$, then $U = F - G$ obeys:
\begin{equation}
D(U) + U = 0
\end{equation}
So $U$ is a decay mode with $\lambda = -1$. This is the operator-level statement that broken symmetry evolves exponentially.

\section{Operator Tower Extraction}

\subsection{Represented Generator}

Let $R: \mathcal{A} \rightarrow \mathcal{V}$ be a representation (finite-dimensional, Banach, symbolic, etc.). Define the induced generator $D_R$ on $R(\mathcal{A})$ by:
\begin{equation}
D_R(R(f)) := R(Df)
\end{equation}
\textbf{Well-definedness condition}: $R \circ D$ factors through the quotient by the collapse ideal $\N$.

\subsection{Operator Tower Levels}

\begin{enumerate}
    \item \textbf{Level 0: base generator} — $D$ (abstract derivation)
    \item \textbf{Level 1: represented generator} — $D_R$ on $R(\mathcal{A})$
    \item \textbf{Level 2: exponential / change operator} — $E_t^R := \exp(tD_R)$ (continuous), $C_{\Delta t}^R := I + \Delta t D_R$ (discrete)
    \item \textbf{Level 3: spectral extractor} — spectrum / eigenmodes / singular values of $D_R$ or $E_t^R$
\end{enumerate}

\subsection{Spectral Invariants}

Assume $D_R$ is a linear operator on a finite-dimensional or spectral-admitting space. Define:
\begin{itemize}
    \item \textbf{Rates}: eigenvalues $\lambda \in \Spec(D_R)$
    \item \textbf{Mode projectors}: $P_\lambda$ such that $D_R P_\lambda = \lambda P_\lambda$
    \item \textbf{Trace invariants}: $\tr(E_t^R)$, $\tr((D_R)^k)$ for integer $k \geq 1$
    \item \textbf{Determinant / zeta invariants}: $\det(E_t^R)$, $\zeta_{D_R}(s)$ (when defined)
\end{itemize}

\subsection{Representation-Change Invariance}

Let $R_1, R_2$ be two representations. A spectral output is \textbf{universal} only if it is invariant under representation change up to conjugacy: there exists an isomorphism $U$ such that $D_{R_2} = U D_{R_1} U^{-1}$. Then $\Spec(D_{R_1}) = \Spec(D_{R_2})$.

\subsection{Collapse Ideal (Invisible Dynamics)}

Define the invisible subspace / collapse ideal for representation $R$:
\begin{equation}
\N_R := \{ f \in \mathcal{A} : R(f) = 0 \}
\end{equation}
If $D(\N_R) \subseteq \N_R$, then invisible dynamics remain invisible. If not, the tower detects "emergence"—invisible components can become visible under evolution.

\section{Noncommuting Flow Composition}

\subsection{Composition of Different Flows}

Let $E_t^{(1)} = \exp(tD_1)$, $E_s^{(2)} = \exp(sD_2)$. Define the composed evolution:
\begin{equation}
E_{t,s} := E_t^{(1)} \circ E_s^{(2)}
\end{equation}
In general, $E_t^{(1)} \circ E_s^{(2)} \neq E_s^{(2)} \circ E_t^{(1)}$. Noncommutativity is measured by the commutator derivation $[D_1, D_2] = D_1 D_2 - D_2 D_1$.

\subsection{Lie-Closure Requirement (BCH Admissibility)}

The set of generators is closed under commutators \textbf{up to a chosen ideal} $\N$ (collapse/gauge ideal):
\begin{equation}
[D_i, D_j] \in \text{Span}\{D_k\} + \N
\end{equation}
Then the composed flow can be represented (on the quotient) as a single exponential:
\begin{equation}
E_t^{(1)} \circ E_s^{(2)} = \exp(\text{BCH}(tD_1, sD_2)) \mod \N
\end{equation}

\subsection{BCH Expansion}

When commutators are meaningful on the domain:
\begin{align}
\text{BCH}(tD_1, sD_2) &= tD_1 + sD_2 + \frac{1}{2} ts [D_1, D_2] \\
&\quad + \frac{1}{12} t^2 s [D_1, [D_1, D_2]] - \frac{1}{12} t s^2 [D_2, [D_1, D_2]] + \cdots
\end{align}
This defines the \textbf{effective generator} of composed change.

\subsection{Integrability Criterion}

If $[D_1, D_2] = 0$, then:
\begin{equation}
E_t^{(1)} \circ E_s^{(2)} = E_s^{(2)} \circ E_t^{(1)} = \exp(tD_1 + sD_2)
\end{equation}
So the multi-flow system reduces to an abelian semigroup.

\subsection{Collapse-Stable Composition}

Let $P: \mathcal{A} \rightarrow \mathcal{A}$ be a collapse / coarse-graining projection. \textbf{Collapse-stability condition}:
\begin{equation}
P \circ E_t = E_t \circ P
\end{equation}
Equivalent generator condition: $P \circ D = D \circ P$. Then evolution descends consistently to the quotient ("evolve then collapse = collapse then evolve").

\section{Discrete Composition and Holonomy}

\subsection{Discrete Step Operators}

Define discrete step operators (Euler truncation):
\begin{equation}
C_1 := I + \Delta D_1, \quad C_2 := I + \Delta D_2
\end{equation}

\subsection{Discrete Composition (Order Dependence)}

Compute:
\begin{align}
C_1 C_2 &= I + \Delta(D_1 + D_2) + \Delta^2 D_1 D_2 \\
C_2 C_1 &= I + \Delta(D_1 + D_2) + \Delta^2 D_2 D_1
\end{align}
Thus:
\begin{equation}
C_1 C_2 - C_2 C_1 = \Delta^2 [D_1, D_2]
\end{equation}
Curvature appears at \textbf{second order} in discrete composition.

\subsection{Discrete Curvature Operator}

Define discrete curvature at step $\Delta$:
\begin{equation}
\mathcal{K}_\Delta(D_1, D_2) := \frac{1}{\Delta^2} (C_1 C_2 - C_2 C_1)
\end{equation}
Then (for Euler steps):
\begin{equation}
\mathcal{K}_\Delta(D_1, D_2) = [D_1, D_2]
\end{equation}
whenever both compositions are defined on the domain.

\subsection{Discrete Commutator Loop (Holonomy Element)}

When inverses exist, define the discrete commutator loop:
\begin{equation}
\mathcal{H}_\Delta := C_1 C_2 C_1^{-1} C_2^{-1}
\end{equation}
For small $\Delta$, the loop expands as:
\begin{equation}
\mathcal{H}_\Delta = I + \Delta^2 [D_1, D_2] + O(\Delta^3)
\end{equation}
This is the \textbf{discrete holonomy} generated by noncommuting change.

\subsection{Finite-Step BCH Effective Generator}

We seek an effective generator $D_{\text{eff}}$ such that $C_1 C_2 \approx \exp(\Delta D_{\text{eff}})$. Consistency forces:
\begin{equation}
D_{\text{eff}} = D_1 + D_2 + \frac{\Delta}{2} [D_1, D_2] + O(\Delta^2)
\end{equation}
So discrete composition induces an extra \textbf{commutator drift} of size $O(\Delta)$ in the effective generator.

\subsection{Symmetric Composition (Commutator Cancellation)}

Define the symmetric (Strang-type) step:
\begin{equation}
C_{\text{sym}} := C_1^{1/2} C_2 C_1^{1/2}
\end{equation}
Then the effective generator satisfies $D_{\text{eff,sym}} = D_1 + D_2 + O(\Delta^2)$. The leading commutator correction cancels.

\subsection{Convergence of Discrete Holonomy}

Let total time be $T$. Consider repeated discrete commutator loops with $\Delta = T/n$:
\begin{equation}
(\mathcal{H}_{T/n})^n
\end{equation}
Under standard boundedness assumptions after passing to a representation space, discrete holonomy converges to the continuous commutator flow generated by $[D_1, D_2]$.

\subsection{Accumulated Residue}

For loops with $[D_1, D_2] \neq 0$, repeated reference (iteration) produces a \textbf{non-vanishing accumulated residue}: nonzero commutator ⇒ nontrivial holonomy accumulation under repeated looping.

\newpage

\part{TRACE-PHASE ALGEBRA AND COCYCLES}

\section{Phase Group and Exponentiation}

\subsection{Phase Group}

Fix an abelian group (or abelian monoid) $(\A, +, 0)$ and a \textbf{phase group} $(\mathcal{P}, \cdot, 1)$ together with an \textbf{exponentiation map}:
\begin{equation}
\exp: \A \rightarrow \mathcal{P}
\end{equation}
satisfying $\exp(a + b) = \exp(a) \cdot \exp(b)$, $\exp(0) = 1$.

Define the \textbf{phase} of an operator:
\begin{equation}
\Phi(O) := \exp(\tr(O)) \in \mathcal{P}
\end{equation}

\subsection{Twisted Trace}

Introduce a map $\omega: \Opr \times \Opr \rightarrow \A$ called the \textbf{Λ-cocycle} satisfying:
\begin{equation}
\omega(O_1, O_2) + \omega(O_1 \circ O_2, O_3) = \omega(O_2, O_3) + \omega(O_1, O_2 \circ O_3)
\end{equation}
and gauge-blindness: if any argument lies in $\N$, then $\omega(N, O) = 0 = \omega(O, N)$.

Define the \textbf{twisted trace} $\tr_\omega$ by:
\begin{equation}
\tr_\omega(O_2 \circ O_1) = \tr_\omega(O_2) + \tr_\omega(O_1) + \omega(O_2, O_1)
\end{equation}
Keep overlay additivity exact: $\tr_\omega(O_1 \otimes O_2) = \tr_\omega(O_1) + \tr_\omega(O_2)$.

\subsection{Holonomy as Completed Invariant}

Define \textbf{holonomy} of a composable word $W = O_n \circ \cdots \circ O_1$ by:
\begin{equation}
\Hol(W) := \Phi_\omega(W) := \exp(\tr_\omega(W)) \in \mathcal{P}
\end{equation}

\begin{theorem}[Holonomy Factorization]
For $W = O_n \circ \cdots \circ O_1$:
\begin{multline}
\tr_\omega(W) = \sum_{k=1}^n \tr_\omega(O_k) \\
+ \sum_{k=1}^{n-1} \omega(O_{k+1} \circ \cdots \circ O_n, O_k \circ \cdots \circ O_1)
\end{multline}
\end{theorem}

\subsection{Phase Equivalence}

Two composable words $W, W'$ are \textbf{phase-equivalent} if:
\begin{enumerate}
    \item they act the same on the quotient: $\bar{W} = \bar{W'}$, and
    \item they have the same holonomy: $\Hol(W) = \Hol(W')$
\end{enumerate}
Write $W \sim_{\text{comp}} W'$. This is the precise "completion" notion: same observable action + same hidden phase memory.

\subsection{Coboundaries and Gauge}

Two cocycles $\omega, \omega'$ are \textbf{cohomologous} if there exists a map $\alpha: \Opr \rightarrow \A$ (gauge 1-cochain) such that:
\begin{equation}
\omega'(O_2, O_1) = \omega(O_2, O_1) + \alpha(O_2 \circ O_1) - \alpha(O_2) - \alpha(O_1)
\end{equation}
and $\alpha(N) = 0$ for $N \in \N$.

\begin{theorem}[Gauge-Invariance]
If $\omega$ and $\omega'$ are cohomologous, then the phase-equivalence relation $\sim_{\text{comp}}$ is unchanged.
\end{theorem}

\subsection{Central Extension}

Define the \textbf{central extension} of $\Opr$ by $\A$:
\begin{equation}
\widetilde{\Opr} := \A \times \Opr
\end{equation}
with composition:
\begin{equation}
(a, O_2) \star (b, O_1) := (a + b + \omega(O_2, O_1), O_2 \circ O_1)
\end{equation}

\begin{theorem}
$\star$ is associative iff $\omega$ satisfies the cocycle law.
\end{theorem}

\subsection{Completed Trace Becomes Homomorphism}

Define:
\begin{equation}
\widetilde{\tr}: \widetilde{\Opr} \rightarrow \A, \quad \widetilde{\tr}(a, O) = a + \tr_\omega(O)
\end{equation}

\begin{theorem}
\begin{equation}
\widetilde{\tr}((a, O_2) \star (b, O_1)) = \widetilde{\tr}(a, O_2) + \widetilde{\tr}(b, O_1)
\end{equation}
\end{theorem}

The extension absorbs the defect and restores strict linear behavior at the trace level.

\newpage

\part{SPECTRAL LAYER}

\section{Characters and Spectrum}

\subsection{Characters as Spectral Points}

\begin{definition}[Character]
A \textbf{character} is a map $\chi: \widetilde{\Opr} \rightarrow \mathcal{P}$ such that for all composable elements:
\begin{equation}
\chi(X \star Y) = \chi(X) \cdot \chi(Y), \quad \chi(\widetilde{\id}) = 1
\end{equation}
and \textbf{gauge-blindness}: $\chi(0, N) = 1$ for all $N \in \N$.
\end{definition}

\begin{definition}[Spectrum]
\begin{equation}
\Spec(\widetilde{\Opr}) := \{ \chi \text{ characters on } \widetilde{\Opr} \}
\end{equation}
\end{definition}

\subsection{Canonical Character from Trace}

Define the canonical phase map:
\begin{equation}
\chi_{\tr}(X) := \exp(\widetilde{\tr}(X))
\end{equation}

\begin{theorem}
$\chi_{\tr}$ is a character.
\end{theorem}

Thus at minimum, the calculus has at least one spectral point (the trace-phase).

\section{Modes as Eigen-Characters of Derivations}

Extend derivations to the completed algebra:
\begin{equation}
\widetilde{D}(a, O) := (a, D(O))
\end{equation}

\begin{definition}[Mode / Eigen-character]
A character $\chi \in \Spec(\widetilde{\Opr})$ is a \textbf{$D$-mode with rate $\lambda \in \A$} if for all $X \in \widetilde{\Opr}$ (where defined):
\begin{equation}
\chi(\widetilde{D}(X)) = \exp(\lambda) \cdot \chi(X)
\end{equation}
\end{definition}

\subsection{Zero-Modes and Invariants}

A character $\chi$ is a \textbf{zero-mode} of $D$ if $\lambda = 0$, i.e., $\chi(\widetilde{D}(X)) = \chi(X)$.

\begin{theorem}
If $O$ is $D$-invariant ($D(O) \in \N$), then for every character $\chi$:
\begin{equation}
\chi(0, O) \text{ is constant along the $D$-flow}
\end{equation}
\end{theorem}

Thus "invariants = zero-modes" is now formal.

\subsection{Spectral Gap and Decay Classes}

Pick a distinguished flow $D$ and consider its modes $(\chi, \lambda)$. Define decay order purely algebraically by comparing rates in $\A$, assuming $\A$ has a preorder $\preceq$.

\begin{definition}[Spectral Gap]
A \textbf{gap} exists if there is a minimal positive rate $\lambda_* \succ 0$ such that:
\begin{itemize}
    \item zero-modes have $\lambda = 0$
    \item every nonzero mode satisfies $\lambda \succeq \lambda_*$
\end{itemize}
\end{definition}

\begin{definition}[Decay Class]
Two flows $D, D'$ are in the same \textbf{decay class} if their mode-rate multisets agree up to scaling by an order-automorphism of $\A$.
\end{definition}

\subsection{Cohomology Class as Spectral Obstruction}

\begin{theorem}
If $[\omega] = 0$ (i.e., $\omega$ is a coboundary), then the completed algebra is equivalent (up to gauge) to an untwisted product, and the spectrum collapses to trace-generated characters.
\end{theorem}

\begin{theorem}
If $[\omega] \neq 0$, then there exist characters not equivalent to $\chi_{\tr}$; these form distinct phase sectors.
\end{theorem}

Thus "completion" literally creates new spectral sectors.

\subsection{Spectral Entropy}

Let $\mathcal{S}(D)$ be the set of mode rates $\{\lambda\}$ (with multiplicity if added).

\begin{definition}[Spectral Entropy]
Define \textbf{spectral entropy} as any monotone functional $\mathbb{H}: \mathcal{S}(D) \rightarrow \A$ satisfying:
\begin{itemize}
    \item $\mathbb{H} = 0$ iff only zero-modes exist
    \item $\mathbb{H}$ increases when new nonzero rates appear or gaps shrink
\end{itemize}
A canonical choice: $\mathbb{H}(D) = \sum_{\lambda \neq 0} \lambda \cdot \dim(\lambda)$ when dimensions are defined.
\end{definition}

\subsection{Functional Calculus}

\begin{definition}[Spectral Integration]
For any function $f: \mathcal{P} \rightarrow \mathbb{C}$ (or appropriate target), define:
\begin{equation}
f(\widetilde{D})(X) := \int_{\Spec(\widetilde{\Opr})} f(\chi(\widetilde{D})) \cdot \chi(X) \, d\mu(\chi)
\end{equation}
where $\mu$ is the Plancherel measure if defined.
\end{definition}

\newpage

\part{CLOCK AND CALCULUS RECOVERY}

\section{Why Calculus Must Reappear}

Up to this point we have:
\begin{itemize}
    \item operators with two products $(\circ, \otimes)$
    \item collapse quotient and completion
    \item derivations $D$
    \item spectrum, modes, decay, entropy
\end{itemize}
But \textbf{no time, no numbers, no limits}. Thus calculus must emerge as a \textbf{representation choice}, not an axiom.

\subsection{Clock as a Morphism}

\begin{definition}[Clock]
A \textbf{clock} is a monoid homomorphism:
\begin{equation}
\tau: (\mathbb{T}, +, 0) \rightarrow (\A, +, 0)
\end{equation}
from a parameter monoid $\mathbb{T}$ ("time labels") into the trace/phase monoid $\A$.
\end{definition}

Interpretation: A clock tells us how many trace-units correspond to one step of evolution.

Examples:
\begin{itemize}
    \item $\mathbb{T} = \mathbb{R}_{\geq 0}$ → continuous time
    \item $\mathbb{T} = \mathbb{Z}_{\geq 0}$ → discrete time
    \item $\mathbb{T} = \mathbb{R}_{\geq 0}$ via entropy → entropy-time
    \item $\mathbb{T} = \mathbb{N}^k$ → multi-parameter time
\end{itemize}

\subsection{Representing Derivations as Flows}

Given a derivation $D$ and a clock $\tau$, define the \textbf{time-labeled flow}:
\begin{equation}
\Phi_t^D(O) := \Phi_{n(t)}^D(O)
\end{equation}
where $n(t)$ is chosen so that $\tau(t) = n(t) \cdot \lambda$ for the dominant spectral rate $\lambda$.

\subsection{Derivative Emerges as Representation Limit}

Let $F: \Opr \rightarrow \mathcal{V}$ be a representation into a linear space $\mathcal{V}$ (with topology). Define the \textbf{derivative of $F$ along $D$} by:
\begin{equation}
\frac{d}{dt} F(O) := \lim_{\Delta t \to 0} \frac{F(\Phi_{\Delta t}^D(O)) - F(O)}{\tau(\Delta t)}
\end{equation}
if the limit exists in the chosen representation.

The derivative:
\begin{itemize}
    \item is \textbf{not algebraic}
    \item exists only after choosing: representation, topology, clock
\end{itemize}
Thus calculus is \textbf{optional structure}, not foundational.

\subsection{Integral as Accumulated Trace}

Define the \textbf{integral along a flow} purely algebraically:
\begin{equation}
\int_0^T O \, dt := \bigotimes_{0 \leq t \leq T} \Phi_t^D(O)
\end{equation}
(formal overlay product). Under a linear representation $F$ that respects overlay as addition:
\begin{equation}
F\left( \int_0^T O \, dt \right) = \int_0^T F(\Phi_t^D(O)) \, d\tau(t)
\end{equation}
This recovers the classical integral.

\subsection{Fundamental Theorem of Calculus (Derived)}

\begin{theorem}[FTC as Representation Theorem]
Let $F$ be a representation where:
\begin{itemize}
    \item overlay → addition
    \item derivation → differentiation
    \item collapse ideal → zero
\end{itemize}
Then:
\begin{equation}
\frac{d}{dt} \left( \int_0^t O \, ds \right) = F(O(t))
\end{equation}
and
\begin{equation}
\int_0^T \frac{d}{dt} F(O(t)) \, dt = F(O(T)) - F(O(0))
\end{equation}
\end{theorem}

This is not an axiom; it is a compatibility condition between derivation, overlay, clock, and representation.

\section{Recovery of Known Calculi}

Different choices of representation + clock recover known theories:

\subsection{Classical Calculus}
\begin{itemize}
    \item $\mathcal{V} = C^\infty(\mathbb{R})$
    \item $\tau(t) = t$
    \item $D = \frac{d}{dx}$
\end{itemize}

\subsection{Stochastic Calculus}
\begin{itemize}
    \item $\mathcal{V} = L^2(\Omega)$
    \item $\tau$ = quadratic variation
    \item cocycle $[\omega] \neq 0$ → Itō correction
\end{itemize}

\subsection{Entropy-Time Calculus}
\begin{itemize}
    \item $\A$ ordered by spectral entropy
    \item $\tau(t) = \mathbb{H}(t)$
    \item flow rates encode irreversibility
\end{itemize}

\subsection{Quantum Mechanics}
\begin{itemize}
    \item $\mathcal{V}$ = Hilbert space
    \item $D = \frac{i}{\hbar}[H, \cdot]$
    \item $\tau(t) = t$ (physical time)
    \item $\otimes$ = pointwise multiplication of operators
\end{itemize}

\subsection{Thermodynamics}
\begin{itemize}
    \item $\mathcal{V}$ = space of observables
    \item $D$ = generator of thermodynamic process
    \item $\tau$ = empirical entropy
    \item trace = free energy
\end{itemize}

All are representations of the same Λ-Morphic Algebra.

\section{Foundational Closure Theorem}

\begin{theorem}[Foundational Closure]
Every calculus satisfying:
\begin{itemize}
    \item chain rule
    \item product rule
    \item fundamental theorem
    \item invariance under reparameterization
\end{itemize}
is equivalent to a representation of a Λ-Morphic Algebra with a clock. Conversely, no additional axioms are required to explain calculus.
\end{theorem}

\newpage

\part{WORKED SYSTEMS}

\section{Rewrite Semigroup Flow (Sūtra-Flow)}

\subsection{State Space and Observables}

\textbf{State space:} Alphabet (tokens) $\Sigma$, configuration space (finite strings) $X := \Sigma^*$.

\textbf{Observable algebra:} $\mathcal{A} := \{ f: X \rightarrow \mathbb{R} \}$ with pointwise addition/multiplication.

\subsection{Rewrite Generators (Partial Operators)}

For each rule index $k \in K$, define a partial rewrite operator $R_k: X \rightarrow X$:
\begin{itemize}
    \item if pattern matches with guard satisfied → rewrite
    \item else → identity on that input
\end{itemize}

\subsection{Ordered Semigroup (Precedence)}

Define a precedence order $\succ$ on rules (vipratiṣedha / blocking). Define precedence composition operator $\triangleright$ on rewrites by:
\begin{itemize}
    \item $R_i \triangleright R_j$ equals $R_i$ alone if $i \succ j$ and both could apply in conflict range
    \item otherwise standard composition $R_i \circ R_j$
\end{itemize}
Define the ordered rewrite semigroup $(\mathcal{R}, \triangleright)$ generated by the $R_k$ under $\triangleright$.

\subsection{One-Step Evolution Operator $U$}

Define a deterministic one-step evolution operator $U: X \rightarrow X$ by applying the precedence-resolved rule schedule:
\begin{equation}
U := R_{k_1} \triangleright R_{k_2} \triangleright \cdots \triangleright R_{k_m}
\end{equation}
Then discrete-time flow: $w_{t+1} = U(w_t)$.

\textbf{Fixed points:} $w_*$ is a normal form iff $U(w_*) = w_*$. Define $NF(w)$ = the fixed point reached by iterating $U$ (when termination + confluence hold).

\subsection{Clock Choice}

\textbf{Clock A: step clock} — $g_A(w_t) := t$

\textbf{Clock B: ranking clock (termination rank)} — Choose $\rho: X \rightarrow \mathbb{N}$ such that every strict rewrite decreases $\rho$: $\rho(U(w)) < \rho(w)$ whenever $U(w) \neq w$. Define $g_B(w) := \rho(w)$.

\subsection{Derivative on This System}

For discrete flow, define the $g$-derivative at step $t$:
\begin{equation}
D_{(U,g)} f(w_t) := \frac{f(w_{t+1}) - f(w_t)}{g(w_{t+1}) - g(w_t)}
\end{equation}
with $w_{t+1} = U(w_t)$, requiring $g(w_{t+1}) \neq g(w_t)$ off fixed points.

\subsection{Generator $D$ and Exponential Lift}

Define the discrete generator (difference generator):
\begin{equation}
D := U - I \quad \text{acting on observables by pullback:} \quad (Uf)(w) := f(U(w)), \quad (Df)(w) := (Uf)(w) - f(w)
\end{equation}
Then "exponential change" operator is defined formally by $C(t) := e^{tD}$.

\subsection{Invariants (Zero-Modes)}

Define the \textbf{normal-form invariant}:
\begin{equation}
I(w) := NF(w) \quad \text{as a function into the set of normal forms}
\end{equation}
Equivalently as an observable family $f_\eta(w) = \mathbbm{1}\{NF(w) = \eta\}$ for each normal form $\eta$. Then along the flow: $NF(U(w)) = NF(w)$, so each $f_\eta$ is a conserved mode: $Df_\eta = 0$ (kernel of $D$).

\subsection{Decay Mode Construction}

Define a distance-to-fixed-point functional $V: X \rightarrow \mathbb{R}_{\geq 0}$ such that:
\begin{itemize}
    \item $V(w) = 0$ iff $w$ is a fixed point
    \item $V(U(w)) \leq \alpha V(w)$ for some $0 \leq \alpha < 1$ off fixed points
\end{itemize}
Then $V$ is a decay mode: $V(w_t) \leq \alpha^t V(w_0)$. With $\lambda := -\log(\alpha) > 0$, we have $V(w_t) \leq e^{-\lambda t} V(w_0)$.

\section{Minimal Holonomy Example}

\subsection{Alphabet and Rules}

Alphabet: $\Sigma = \{a, i, e, x\}$

Rules:
\begin{itemize}
    \item Rule L (lopa / deletion): $x \to \varepsilon$
    \item Rule S1 (sandhi-with-marker): $a x i \to e$
    \item Rule S2 (sandhi-without-marker): $a i \to e$
\end{itemize}

\subsection{Base Evolution and Diamond}

Initial word: $w_0 = axi$. Target normal form: $NF(w_0) = e$.

Two derivation paths:
\begin{itemize}
    \item Path P (direct sandhi): $axi \xrightarrow{S1} e$ — P = [S1]
    \item Path Q (lopa then sandhi): $axi \xrightarrow{L} ai \xrightarrow{S2} e$ — Q = [L, S2]
\end{itemize}
Both end at $e$, so base system is confluent at this diamond.

\subsection{Holonomy Data ($\lambda$-Lift)}

Lift to $\widetilde{X} = X \times U(1)$. Assign rule phases:
\begin{equation}
\chi_\lambda(S1) = 1, \quad \chi_\lambda(S2) = 1, \quad \chi_\lambda(L) = \exp(i\lambda)
\end{equation}
Use state-independent cocycle: $\omega_\lambda(k,w) = \chi_\lambda(k)$. Lifted operators:
\begin{equation}
\widetilde{R}_k(w,\theta) = (R_k(w), \chi_\lambda(k)\theta)
\end{equation}

\subsection{Compute Holonomy}

Start with $\theta_0 = 1$.

Path P: apply S1 once → multiplier $1$ → holonomy $H_\lambda(P) = 1$, final state $(e,1)$.

Path Q: apply L then S2 → multiplier $\exp(i\lambda) \cdot 1$ → holonomy $H_\lambda(Q) = \exp(i\lambda)$, final state $(e,\exp(i\lambda))$.

\subsection{Result}

We have constructed:
\begin{itemize}
    \item identical projected outcome: $\pi(\widetilde{w}_P) = \pi(\widetilde{w}_Q) = e$
    \item different trace phase: $\theta_P \neq \theta_Q$ when $\lambda \not\equiv 0 \mod 2\pi$
\end{itemize}
This is the minimal "hidden phase cancels in output but persists in trace" mechanism.

\subsection{Flatness Constraint}

If we require phase path-independence for confluent derivations, we must impose $1 = \exp(i\lambda)$, so $\lambda \equiv 0 \mod 2\pi$.

\subsection{Spectral Consequence}

Define the lifted one-step operator $\widetilde{U}$. Base dynamics has a single absorbing normal form $e$. Lifted dynamics splits into phase sectors when both paths are allowed: one sector with phase 1, one with phase $e^{i\lambda}$.

\section{Torus System (Compact Phase)}

\subsection{State Space and Observable Algebra}

State space: $M := T^n$ (angles $\theta \in (S^1)^n$). Observable algebra: $\mathcal{A} := C^\infty(T^n, \mathbb{C})$.

\subsection{Flow and Generator}

Choose a smooth vector field:
\begin{equation}
v = \sum_{j=1}^n \omega_j(\theta) \frac{\partial}{\partial \theta_j}
\end{equation}
Define generator as Lie derivative along $v$: $D := D_v: \mathcal{A} \rightarrow \mathcal{A}$. Define the \textbf{change step} operator: $C_v := I + D_v$.

\subsection{Change-Invariance Equation → Decay Mode}

Given two observables $F, G \in \mathcal{A}$, impose $C_v(F) = C_v(G)$. Equivalent to $(F-G) + D_v(F-G) = 0$. Define $U := F - G$. Then $D_v(U) = -U$, so $U$ is a decay mode with rate $\lambda = 1$.

\subsection{Explicit Solution Along Flow Lines}

Let $\gamma(t)$ be an integral curve of $v$. Then:
\begin{equation}
\frac{d}{dt} (U \circ \gamma) + (U \circ \gamma) = 0 \Rightarrow U(\gamma(t)) = U(\gamma(0)) e^{-t}
\end{equation}

\subsection{Spectral Basis on $T^n$}

Character basis (Fourier modes): $\chi_k(\theta) = \exp(i\sum_{j=1}^n k_j \theta_j)$, $k \in \mathbb{Z}^n$. Assume constant rotational field: $v = \sum_{j=1}^n \Omega_j \frac{\partial}{\partial \theta_j}$. Then $D_v \chi_k = i(k \cdot \Omega) \chi_k$.

\subsection{Compatibility / Obstruction}

Expand $U = \sum_k \widehat{U}_k \chi_k$. Insert into $U + D_v U = 0$:
\begin{equation}
(1 + i k \cdot \Omega) \widehat{U}_k = 0 \quad \forall k
\end{equation}
For real $\Omega$, the condition $k \cdot \Omega = i(-1)$ is impossible, so all nonzero modes vanish. Only the zero-mode survives.

\subsection{Remedy Constructions}

\textbf{Clock change:} Choose clock $g$ along flow lines, define $D^{(g)} := \frac{1}{\dot{g}} D$. Decay equation becomes $D(U) + \dot{g} U = 0$.

\textbf{Step operator change:} Replace $I + \Delta D$ by implicit step $S_{\Delta}^{\text{imp}} := (I - \Delta D)^{-1}$ or Cayley step $S_{\Delta}^{\text{cay}} := (I - \frac{\Delta}{2} D)^{-1} (I + \frac{\Delta}{2} D)$.

\textbf{Sectoring by constraints:} Restrict to constrained subspace $T_\Gamma^n := \{\theta \in T^n : A\theta = c \ (\text{mod } 2\pi)\}$.

\section{Probabilistic / Stochastic System}

\subsection{State Space}

Let $X = \{1,2,\ldots,N\}$ (finite state space). Probability simplex:
\begin{equation}
\mathcal{P}(X) := \left\{ p \in \mathbb{R}^N \mid p_i \geq 0, \sum_i p_i = 1 \right\}
\end{equation}
Observables: $\mathcal{A} := \{ f: X \rightarrow \mathbb{R} \} \cong \mathbb{R}^N$.

\subsection{Stochastic Evolution (Markov Operator)}

Let $K$ be a stochastic matrix: $K_{ij} \geq 0$, $\sum_j K_{ij} = 1$. State evolution: $p_{n+1} = p_n K$. Observable evolution (Koopman): $(Tf)(i) = \sum_j K_{ij} f(j)$.

\subsection{Discrete Generator}

Define $D := T - I$, so $(Df)(i) = \sum_j K_{ij} f(j) - f(i)$.

\subsection{Continuous-Time Limit}

Assume $K = I + \Delta L + o(\Delta)$ where $L_{ij} \geq 0$ ($i \neq j$), $\sum_j L_{ij} = 0$. Then as $\Delta \to 0$: $\frac{d}{dt} f_t = L f_t$, generator $D = L$.

\subsection{Entropy Clock}

Define entropy $S(p) := -\sum_i p_i \log p_i$, clock $g(t) := S(p_t)$. This is monotone for irreducible Markov chains. Clocked derivative $D_g f := \frac{Df}{Dg}$.

\subsection{Spectrum (Rates)}

Solve $L u_k = \lambda_k u_k$:
\begin{itemize}
    \item $\lambda_0 = 0$ (stationary mode)
    \item $\Re(\lambda_k) < 0$ for $k \geq 1$
\end{itemize}
Evolution: $f_t = \sum_k e^{\lambda_k t} c_k u_k$. These $\lambda_k$ are decay rates / relaxation times.

\subsection{Invariants (Kernel Modes)}

Stationary distribution $\pi$ satisfies $\pi L = 0$. Invariant observable: $Df = 0 \iff f \in \text{Span}\{\mathbf{1}\}$. Kernel: $\ker(D) = \text{stationary subspace}$.

\subsection{Collapse / Śūnya}

Define coarse-graining projection $(Pf)(i) := \sum_{j \in C(i)} \frac{1}{|C(i)|} f(j)$. If $PL = LP$, then coarse-grained evolution is Markov, entropy clock collapses micro-distinctions, kernel enlarges.

\section{Physical System (Diffusion)}

\subsection{State Space}

Let $M$ be a domain (manifold or subset of $\mathbb{R}^d$) with boundary conditions. Probability densities:
\begin{equation}
\mathcal{P}(M) = \{\rho \geq 0, \int_M \rho \, dx = 1\}
\end{equation}
Observables: $\mathcal{A} = \{ f: M \rightarrow \mathbb{R} \}$.

\subsection{Physical Flow (Diffusion)}

Density evolution: $\partial_t \rho = \nabla \cdot (\kappa(x) \nabla \rho)$ where $\kappa(x) \geq 0$ is diffusivity. Observable evolution (adjoint): $\partial_t f = \kappa(x) \Delta f + \nabla \kappa(x) \cdot \nabla f$.

\subsection{Generator}

Define generator $L$ acting on observables: $Lf := \nabla \cdot (\kappa \nabla f)$. Evolution semigroup: $E_t := e^{tL}$, so $f_t = E_t f_0$.

\subsection{Heat Kernel Representation}

There exists a kernel $K_t(x,y)$ such that $(E_t f)(x) = \int_M K_t(x,y) f(y) \, dy$.

\subsection{Spectrum (Rates) and Mode Law}

Assume $L$ is self-adjoint. Solve $L \phi_n = -\lambda_n \phi_n$, $\lambda_n \geq 0$. Then $E_t \phi_n = e^{-\lambda_n t} \phi_n$. Spectral expansion:
\begin{equation}
f_t = \sum_n e^{-\lambda_n t} \langle f_0, \phi_n \rangle \phi_n
\end{equation}

\subsection{Invariants (Kernel Modes)}

$\ker(L) = \text{Span}\{1\}$ on connected domain with Neumann/periodic boundaries. Mass conservation: $\frac{d}{dt} \int_M \rho_t \, dx = 0$.

\subsection{Collapse / Śūnya}

Coarse-graining projection $(Pf)(x) := \frac{1}{|B(x)|} \int_{B(x)} f(z) \, dz$. Collapse-stability condition: $PL = LP$ ensures evolution descends consistently.

\section{Learning / Representation System}

\subsection{State Space}

Parameter space $\Theta \subseteq \mathbb{R}^p$, data distribution $x \sim \mathcal{D}$, model $f_\theta(x) \in \mathbb{R}$. Loss functional:
\begin{equation}
\mathcal{L}(\theta) := \mathbb{E}_{x \sim \mathcal{D}} [\ell(f_\theta(x), y(x))]
\end{equation}

\subsection{Morphic Flow (Learning Dynamics)}

Continuous-time gradient flow: $\dot{\theta}_t = -\nabla_\theta \mathcal{L}(\theta_t)$.

\subsection{Observable Algebra}

Observables as real-valued functionals: $\mathcal{A} := \{ F: \Theta \rightarrow \mathbb{R} \}$. Generator:
\begin{equation}
(DF)(\theta) := \nabla_\theta F(\theta) \cdot \dot{\theta} = -\nabla_\theta F(\theta) \cdot \nabla_\theta \mathcal{L}(\theta)
\end{equation}

\subsection{Representation (Feature Lift)}

Representation map $R(\theta) := \phi_\theta \in \mathcal{H}$ where $\mathcal{H}$ is a linear function space over inputs. Canonical choice: $\phi_\theta(x) := f_\theta(x)$. Indistinguishability: $\theta_1 \sim_R \theta_2$ iff $f_{\theta_1}(x) = f_{\theta_2}(x)$ a.e. $x \sim \mathcal{D}$.

\subsection{Induced Generator on Function Space}

Let $u_t(x) := f_{\theta_t}(x)$. Then:
\begin{equation}
\partial_t u_t(x) = \nabla_\theta f_{\theta_t}(x) \cdot \dot{\theta}_t = -\nabla_\theta f_{\theta_t}(x) \cdot \nabla_\theta \mathcal{L}(\theta_t)
\end{equation}

\subsection{Kernel Operator}

Define the kernel $K_\theta(x,x') := \nabla_\theta f_\theta(x) \cdot \nabla_\theta f_\theta(x')$. For squared loss:
\begin{equation}
\partial_t u_t(x) = -\int K_{\theta_t}(x,x') (u_t(x') - y(x')) \, d\mathcal{D}(x')
\end{equation}

\subsection{Linearization Regime}

If kernel is approximately constant: $K_{\theta_t} \approx K_{\theta_0} =: K$, then $\partial_t (u_t - y) = -K(u_t - y)$, so $u_t - y = e^{-tK}(u_0 - y)$. Generator in representation space: $D_R = -K$.

\subsection{Spectrum (Learning Rates)}

Solve $K \psi_n = \lambda_n \psi_n$, $\lambda_n \geq 0$. Then mode law:
\begin{equation}
\langle u_t - y, \psi_n \rangle = e^{-\lambda_n t} \langle u_0 - y, \psi_n \rangle
\end{equation}
$\lambda_n$ are \textbf{learning rates} (per-mode decay rates).

\subsection{Invariants}

Modes with $\lambda_n = 0$ satisfy $\langle u_t - y, \psi_n \rangle = \langle u_0 - y, \psi_n \rangle$. So $\ker(K) = \text{unlearnable subspace}$.

\subsection{Collapse via Compression}

Define a compression projection $P: \mathcal{H} \rightarrow \mathcal{H}$ (low-rank projection, pruning, quantization). Collapse-stability condition: $PK = KP$. Then training commutes with compression: $Pe^{-tK} = e^{-tK}P$.

\newpage

\part{PROGRAM MANIFEST AND CONCLUSIONS}

\section{Program Manifest}

This work develops a framework in which recognition, not computation, is treated as the primitive operation underlying mathematics, physics, and computation.

Classical formalisms begin by constructing objects and defining procedures upon them. In contrast, this framework begins with the observation that \emph{identity must already be recognized} before any construction, algorithm, or inference can occur.

From this standpoint:
\begin{itemize}
    \item \textbf{Computation} arises when recognition cannot be asserted directly.
    \item \textbf{Coherence} measures the stability of recognition under evolution.
    \item \textbf{Partial differential equations} encode local coherence constraints.
    \item \textbf{Spectral objects} arise as residues of constrained coherence.
    \item \textbf{Representations and symmetries} formalize how recognition survives under transformation.
\end{itemize}

Breakdown phenomena—turbulence, undecidability, singularities, spectral instabilities—are not treated as anomalies. They are interpreted as structural failures of recognition under evolution.

This framework does not propose new axioms to solve existing problems prematurely. Its purpose is to reposition problems correctly, identifying where computation is being applied in place of recognition.

The same structural obstruction appears across disparate domains:
\begin{itemize}
    \item nonlinear PDEs,
    \item spectral geometry,
    \item number theory,
    \item quantum theory,
    \item complex systems.
\end{itemize}

This recurrence suggests a common underlying mechanism rather than isolated technical difficulty.

The aim of this program is not immediate resolution, but clarity of structure. Only when problems are situated at the correct structural level can solutions—if they exist—be meaningfully pursued.

\section{Final Synthesis}

We have established:

\begin{enumerate}
    \item \textbf{Recognition as primitive} — $\R(\mathbf{A},\mathbf{B})$ as unmediated sameness.
    \item \textbf{Coherence as recognition preservation} — $\R(\mathbf{A},\mathbf{B}) \Rightarrow \R(\mathbf{E}_t(\mathbf{A}),\mathbf{E}_t(\mathbf{B}))$.
    \item \textbf{Computation as compensatory} — arising only where recognition fails.
    \item \textbf{Calculus as representation choice} — derivative with respect to arbitrary clock/measure $g$.
    \item \textbf{Chain, product, quotient rules} — holding without coordinates.
    \item \textbf{Fundamental theorem} — integration as inverse of differentiation under chosen clock.
    \item \textbf{Higher-order structure} — commutators, curvature, non-integrability.
    \item \textbf{Exponential change operator} — $E_t = \exp(tD)$, discrete truncation $C_{\Delta t} = I + \Delta t D$.
    \item \textbf{Spectral invariants} — modes, zero-modes, decay rates, operator towers.
    \item \textbf{Trace-phase algebra} — cocycles, holonomy, central extensions.
    \item \textbf{Four-layer structure} — Śūnya, Cut, Morphisum, Operator.
    \item \textbf{Worked systems} — rewrite semigroups, holonomy diamonds, torus systems, stochastic processes, diffusion, learning dynamics.
\end{enumerate}

\section{The Śūnya Limit}

The Śūnya limit $\overset{\varnothing}{\lim}_{t \to \infty} x_t := \mathcal{P}(x_t)$ replaces classical pointwise convergence with \textbf{collapse into indistinguishability}. This is the proper mathematical formulation of:
\begin{itemize}
    \item regime change
    \item compression
    \item collapse
    \item renormalization
    \item phase transitions
\end{itemize}

\section{Universality Theorem (Final Form)}

\begin{theorem}[Complete Universality]
Any system satisfying:
\begin{itemize}
    \item a semigroup evolution $\Phi_t: \T \rightarrow \T$
    \item a monotone distinguishability measure $\tau: \mathbb{T} \rightarrow \A$
    \item a representation allowing linearization $F: \Opr \rightarrow \Hom(V)$
\end{itemize}
admits:
\begin{itemize}
    \item a derivative $Df = \lim_{\Delta t \to 0} \frac{f \circ \Phi_{\Delta t} - f}{\tau(\Delta t)}$
    \item an integral $\int f \, d\tau = \lim_{n \to \infty} \bigotimes_{i=1}^n f \circ \Phi_{t_i}$
    \item exponential evolution $U_t = e^{tD}$
    \item spectral invariants $\sigma(D)$
    \item a meaningful limit $\lim_{t \to \infty} f \circ \Phi_t = \Pi_{\ker D}(f)$
\end{itemize}
All classical calculi arise as special cases under choice of:
\begin{itemize}
    \item state space $\T$
    \item clock/measure $\tau$
    \item representation $F$
\end{itemize}
\end{theorem}

\section{Closing Remarks}

This framework does not claim to describe reality. It provides a functional scaffolding within which any consistent mathematical description can appear, function, and dissolve when conditions change. It is, in the deepest sense, a \textit{morphic} foundation—one that takes its own provisionality as its only constant.

\begin{center}
\fbox{
\begin{minipage}{0.8\textwidth}
\centering
\textbf{All constructions are empty.}\\
\textbf{Emptiness itself is empty.}\\
\textbf{Use is allowed.}\\
\textbf{Reification is not.}
\end{minipage}
}
\end{center}

\begin{theorem}[Self-Consistency]
The framework presented here satisfies its own criteria: it is conditional, functional, and collapsible. No axiom claims ultimacy. No construction reifies itself. The system is closed under its own emptiness guard.
\end{theorem}

\begin{proof}
By the emptiness guard, if any reader finds contradiction, the construction dissolves. No repair is required. The proof is the absence of need for proof.
\end{proof}

\end{document}